\chapter{Cartasis}
\label{chap:cartasis}

\section{The bounce membrane}

\Cartasis is the 3D hypersurface in spacetime where torsion-mediated repulsion
equals gravitational attraction, and across which the geometry transitions from
collapsing-parent to expanding-baby. The name captures both the
surface-as-membrane and the moment-as-genesis aspects: it is simultaneously a
spatial boundary in supraverse coordinates and a temporal beginning in
baby-universe coordinates.

\section{Physics of the bounce}

\subsection{The competing terms}

In Einstein--Cartan with fermionic matter, the effective stress--energy at high
density has two relevant contributions. The \emph{gravitational term}
(attractive at all densities),
\[
  T_{\text{grav}} \sim \rho c^2,
\]
scales linearly with mass--energy density. The \emph{torsion term} (repulsive at
high spin density),
\[
  T_{\text{tors}} \sim \frac{G \hbar^2 s^2}{c^4},
\]
scales as the square of the spin density $s$.

At ordinary densities, $T_{\text{tors}} / T_{\text{grav}} \sim 10^{-40}$,
negligible. As matter compresses, both terms grow, but spin density grows faster
(because compression packs more fermions into each volume, each contributing
$\hbar/2$). The quadratic scaling means $T_{\text{tors}}$ overtakes
$T_{\text{grav}}$ at sufficient compression.

\subsection{The critical density}

Setting $T_{\text{grav}} = T_{\text{tors}}$ and solving for density,
\[
  \rhoc \sim \frac{c^4}{G \hbar^2 (s/\rho)^2}.
\]
For typical fermionic matter where $s/\rho \sim \hbar/m_n$ (one fermion-spin per
nucleon-mass), this gives
\[
  \rhoc \sim \frac{m_n^2 c^4}{G \hbar^4} \sim 10^{50}\,\mathrm{kg\,m^{-3}}.
\]
This is 33 orders of magnitude beyond nuclear density and 46 orders of magnitude
below Planck density. The exact value depends on the fermion content (different
masses give different ratios) and on detailed coefficients in the
Einstein--Cartan Lagrangian.

\subsection{What is actually at Cartasis density}

At $\rhoc$, conventional particle descriptions break down. The matter is
compressed beyond quark deconfinement, beyond electroweak symmetry restoration,
beyond any regime where the Standard Model's standard quasiparticle excitations
are well-defined. What exists is \textbf{continuous field content}:
energy--momentum density, spin density, chiral current density. Individual
particle identities do not survive.

What crosses Cartasis is therefore not ``particles'' but \textbf{conserved
currents}: total energy--momentum, total spin, total chiral charge, and
(subject to electroweak anomaly considerations) baryon and lepton numbers.

This has implications for chirality dynamics (see Chapter~\ref{chap:chirality})
and for what kind of structures can be transmitted through a bounce (none---
structures must be rebuilt from elementary field content on the other side).

\section{Geometry of the bounce}

\subsection{The radial--temporal exchange}

Inside a Schwarzschild black hole's event horizon, the radial coordinate $r$
becomes timelike. Future-directed motion is radially inward. This is standard
general relativity.

At Cartasis density, the bounce dynamics reverse this: the torsion repulsion
pushes matter outward in $r$, but since $r$ is timelike inside the horizon,
``outward in $r$'' does not take matter back through the horizon---it takes
matter into the \emph{future} of the bounce.

The future of the bounce, in baby-universe coordinates, is the cosmological
future of the new spacetime region. The radial direction $r$ on the parent's
side smoothly transitions into the time direction $\tau$ on the baby's side at
the bounce hypersurface.

\subsection{Junction conditions}

The metric and connection must be continuous across the bounce. Specifically:

\begin{itemize}
  \item the induced metric on the Cartasis hypersurface matches from both sides;
  \item the extrinsic curvature matches (modified Israel--Darmois conditions
    accommodating torsion);
  \item stress--energy and spin current flow continuously through the surface.
\end{itemize}

These junction conditions are the technical specification of what ``the bounce
is continuous'' means. They are a system of equations relating parent-side and
baby-side field values at the membrane.

\subsection{Bounce as persistent interface, not single event}

Critical point: in a black hole that continues to accrete matter from its parent
universe's environment, the bounce is \textbf{ongoing}. New matter falls in,
compresses, reaches Cartasis density, and joins the bounce. The Cartasis
hypersurface is therefore not a single moment but a persistent interface that
exists for as long as the parent black hole exists.

From the baby universe's interior, this means matter continues to be injected
across the Cartasis membrane throughout the parent's lifetime---see
Chapter~\ref{chap:darksector} for the dark-matter interpretation this enables.

\section{The geometric identification: parent's horizon $=$ baby's cosmic
  horizon}

From the parent's exterior, Cartasis sits inside the event horizon. From the
parent's interior, it is the surface where infalling matter is currently
reaching critical density. From the baby's interior, looking backward in time,
Cartasis is \textbf{the boundary of the past at maximum cosmic lookback}.

The Big Bang membrane (from the baby's perspective) and the parent's
event-horizon region (from the parent's perspective) are the same surface,
viewed from different sides. This is the geometric basis for the
CMB-as-Hawking-radiation identification (see Chapter~\ref{chap:darksector}).

\section{Calculation targets}

To make Cartasis quantitative, the following calculations are needed.

\begin{enumerate}
  \item \textbf{Precise $\rhoc$ from the Einstein--Cartan Lagrangian.} Existing
    estimates vary by factors of 10--100 depending on coefficient choices. A
    clean calculation would pin the value.
  \item \textbf{Bounce timescale.} How long, in proper time, does the bounce
    process take? This sets the timescale for the baby's ``Big Bang'' phase.
  \item \textbf{Bounce efficiency.} What fraction of the parent's infalling
    mass--energy is transmitted to the baby's expanding phase, vs.\ retained in
    some intermediate region, vs.\ scattered back?
  \item \textbf{Chirality selectivity.} Does the bounce treat the two
    chiralities symmetrically (no filter), with mild preference (weak filter),
    or with strong preference (strong filter)? See Chapter~\ref{chap:chirality}.
  \item \textbf{Spin coherence requirements.} Does the bounce require correlated
    spins, or does randomized spin density suffice? Existing analysis suggests
    random spins work because the $s^2$ scaling is positive regardless of
    orientation, but this needs verification.
\end{enumerate}

\section{Open questions specific to Cartasis}

\begin{itemize}
  \item \textbf{Is $\rhoc$ universal?} Or does it depend on the matter content
    (different fermion species, different bound-state effects)?
  \item \textbf{What happens at the boundary of Cartasis density}---is there a
    sharp transition or a gradient region?
  \item \textbf{How does Cartasis interact with rotation?} Kerr-like bouncing
    black holes vs.\ Schwarzschild-like.
  \item \textbf{Does Cartasis preserve baryon number?} Or does the electroweak
    anomaly at trans-EW temperatures violate $B$ explicitly during the bounce?
  \item \textbf{What is the relationship between Cartasis and quantum gravity
    proper?} Cartasis is in a regime where Einstein--Cartan is presumably still
    classical, but full quantum-gravity effects might modify the picture at
    higher densities.
\end{itemize}
